% ============================================
% LUC-13 (EN)
% English translation — SPEC-935-R355
% ============================================
\section*{\textbf{LUC-13} --- The Night of the Diagnosis}
\markboth{LUC-13 --- The Night of the Diagnosis}{}

\noindent \textit{Document extracted from the expert file.}

\vspace{0.5em}\noindent\rule{\textwidth}{0.5pt}\vspace{0.5em}


\noindent \textbf{April 2026 --- Belo Horizonte}\\
\noindent \textbf{Audio recording --- file L260-13A.wav}\\
\noindent \textbf{Duration:} 47 minutes and 32 seconds.\\
\noindent \textbf{Status:} not transcribed. Dr. Lúcia Mendes did not authorize the transcription.

\noindent \textit{Archivist's note: this audio file was found on Dr. Lúcia Mendes's computer, dated April 12, 2026. It is her first personal record about the case. The recording was made at home, at 3:47 in the morning. We transcribed the first 4 minutes and 12 seconds --- the rest is unintelligible.}
\noindent \textbf{[00:00:00] Lúcia:} Test. Audio test. L260... L260... thirteen. A. Recording.

\noindent \textbf{[00:00:12]\textit{ It is the third night. Third night I wake up at the same time. 3:14. Always 3:14. No use adjusting the clock. No use not looking at the time. I }know} it is 3:14 before I look. I feel it in my breathing. In the weight of the air.

\noindent \textbf{[00:00:41]} I have tried every hypothesis. Anticipatory anxiety. Classical conditioning. Post-hypnotic suggestion induced by reading. Everything I would teach in a psychopathology course for residents. Fifteen years of clinical practice and I know exactly how a patient describes a recurrent nocturnal awakening and how one differentiates a sleep disorder from persecutory ideation.

\noindent \textbf{[00:01:18]} But it is not a patient. It is me.

\noindent \textbf{[00:01:25]\textit{ And what I feel when I wake up is not anxiety. It is }expectation}. As if something were going to happen. As if something were \textit{getting ready} to happen. The body waits. The body \textit{knows} something the mind has not yet authorized.

\noindent \textbf{[00:02:03]} The smell came back today. I had spent the whole day without feeling it. I thought: it passed. It was suggestion. A tired brain fabricates stimuli where none exist. Then I turned off the bedroom light, lay down, closed my eyes... and the smell was on the pillow. Sweet. Not perfume-sweet. Sweet of vegetable matter wilting in standing water. Sweet of a bouquet forgotten in a dirty glass until the water turned black.

\noindent \textbf{[00:02:47]} I got up, changed the pillowcase. Smelled the new pillow. Clean. Lavender. The smell continued. It was not on the pillow. It was inside my nose.

\noindent \textbf{[00:03:12]} I turned on every light. Opened the window. Breathed deeply. The street air was cold, normal, urban. Pollution. Diesel. Wet asphalt. Real things. I closed the window. The smell came back at the same intensity.

\noindent \textbf{[00:03:41]} It is \textit{in me}. Not in the environment. In the mucosa. In the back of the throat. In the place where olfactory memory lodges and never leaves again. Like the smell of one's childhood home. Like the smell of someone who is gone.

\noindent \textbf{[00:04:12]} I went to the bathroom and looked in the mirror. The right eye...

\noindent \textit{Here the recording is interrupted for 23 seconds. What follows was partially recovered through audio restoration. Lúcia's voice is different: lower, slower, as if speaking demanded physical effort.}
\noindent \textbf{[00:04:35] Lúcia:} ...is lighter. Not much. Enough. The iris lost pigment on the outer edge. It became the color of... I do not know. It is not yellow yet. It is a shade that has no name. A shade I do not want to name.

\noindent \textbf{[00:05:11]} I touched the eye with my fingertip. The cornea was warm. Eyes are not warm. Eyes are cold. The temperature of the ocular surface is 32 degrees. What I felt was more. Much more.

\noindent \textbf{[00:05:44]\textit{ I made a diary. No. I made a }record}. A spreadsheet. Symptoms, times, intensity, frequency. It lasted three days. On the third day, I realized the spreadsheet was identical to the table Joaquim describes in his diary. The same structure. The same columns. The same words. My handwriting, but the same structure.

\noindent \textbf{[00:06:28]\textit{ I did not copy it. I swear I did not copy it. The table }came}. It came ready. The hand wrote by itself.

\noindent \textbf{[00:07:01]} I tore up the spreadsheet. Threw it in the trash. Then I took it out of the trash, burned it in the sink, washed away the ashes. I do not want anyone to find it.

\noindent \textbf{[00:07:33]} If I read an account like this coming from a patient, my diagnosis would be: paranoid ideation with somatization, possible mild dissociative picture, post-traumatic stress triggered by contact with traumatic material.

\noindent \textbf{[00:08:12]} But what I feel is not paranoia. Paranoia is thinking someone is watching you. What I feel is different. I feel that someone \textit{is already inside}. What Joaquim describes as \textquotedblleft{}the weight behind the sternum\textquotedblright{}, \textquotedblleft{}the breathing that is not mine\textquotedblright{}, \textquotedblleft{}the mouth that moves without my telling it to\textquotedblright{}. I know what he meant. Because I am feeling it.

\noindent \textbf{[00:08:52]} In the early hours I sat down to write. Something I wanted to record. I picked up the notebook. Opened to the first page. The hand wrote: \textit{there is so much hunger here, Lúcia. So much hunger I do not even know where to fit.}

\noindent \textbf{[00:09:07]\textit{ That sentence is in the diary. It is the entity's first speech. I }read} it. But when the hand wrote it, the handwriting was not mine. It was the same handwriting Dr. Oliveira described in his letter: slanted to the left, open \textquotedblleft{}a\textquotedblright{}, irregular pressure. Joaquim's handwriting.

\noindent \textbf{[00:09:41]\textit{ I did not copy it. I did not }need} to copy it. The hand knew the way by itself. The hand remembered a pattern I had seen on paper and reproduced it without authorization. The hand has its own memory.

\noindent \textbf{[00:10:03]} The same pattern. The same structure. Me writing Joaquim's sentence. Me feeling what he felt. Me seeing what he saw. Me.

\noindent \textbf{[00:10:28]} Before, I said \textquotedblleft{}Joaquim describes\textquotedblright{}, \textquotedblleft{}Joaquim reports\textquotedblright{}, \textquotedblleft{}the patient says\textquotedblright{}. Now I say \textquotedblleft{}the trench\textquotedblright{}, \textquotedblleft{}the eye\textquotedblright{}, \textquotedblleft{}she\textquotedblright{}. His proper name is disappearing from my vocabulary. He is no longer the patient. He is the witness. The witness who testifies inside me.

\noindent \textbf{[00:11:12]} Pretty soon I will say \textquotedblleft{}we\textquotedblright{}.

\noindent \textit{[The recording captures a long noise, similar to pages being turned quickly.]}
\noindent \textbf{[00:12:04]} It is not coming. It is not coming. It is not coming. She cannot be here. She cannot be inside me. I am a psychiatrist. I am 42 years old. I do not believe in possession. I do not believe in contagion. I do not believe in anything, anything, anything...\margnota{the sound of\\pages being\\turned.\\Who is\\turning the\\pages?}

\noindent From this point on, the recording deteriorates. Lúcia's voice alternates between long silences and short, disconnected sentences. The audio quality worsens --- there is a background noise that the technicians describe as \textquotedblleft{\textit{similar to wind passing through a narrow duct\textquotedblright{}.}}

\noindent \textbf{[00:12:04] Lúcia:} The diary is in the office. I locked the door.

\noindent \textbf{[00:12:31]} It is useless. I feel it on the other side of the wall. The cover. The weight. The spine.

\noindent \textbf{[00:13:02]\textit{ It }radiates}.

\noindent \textbf{[00:13:28]} At an hour like this... 3:47 in the morning... and I know it is awake. The diary is awake. The pages are moving by themselves in the dark. The ink is running to places I did not write.

\noindent \textbf{[00:14:15]} This is not scientific. This is not clinical. This is the account of an adult woman, with a postdoctoral degree, whose brain is lying to her.

\noindent \textbf{[00:14:52]} \textit{(long pause)} And what if it is not a lie?

\noindent The following 33 minutes consist mostly of silence, irregular breathing, and background noise. At 47 minutes and 12 seconds, the recording ends with what appears to be Lúcia's voice whispering: \textquotedblleft{\textit{Patient 1,262.\textquotedblright{} The recording was manually ended at 47:32. No other audio file was found on the computer.}

\noindent Archivist's note: the morning after this recording, Lúcia Mendes called the Public Archive of Minas Gerais and requested the return of the diary. She said she needed to \textquotedblleft{\textit{confirm some information}\textquotedblright{}} on the back cover. The diary was picked up by her in person the same day. She did not return it until May 13.}




\vspace{0.5em}
\noindent\textit{\textbf{↪ Routes:} \rota{DOC-01} \rotasep \rota{LUC-02} \rotasep \rota{EPI-01} \rotasep \rota{LUC-14}}